\section*{Capítulo \ref{ch:g4:measure} --- Medidas y perímetro}

\begin{solution}{exo:g4:measure:1}
$7$ metro $= 700$ centímetro; \quad $3.2$ kilómetros $= 3\,200$ metro; \quad
$85$ milímetros $= 8.5$ centímetro; \quad $4\,500$ metro $= 4.5$ km.
\end{solution}

\begin{solution}{exo:g4:measure:2}
$2$ kilos $= 2\,000$ gramo; \quad $1\,800$ gramo $= 1$ kilos $800$ gramo; \quad
$3.5$ l $= 350$ CL; \quad $25$ CL $= 0.25$ l.
\end{solution}

\begin{solution}{exo:g4:measure:3}
Rectángulo: $2 \times (9 + 5) = 28$ centímetro. Cuadrado:
$4 \times 7.5 = 30$ centímetro. Triángulo: $6 + 8 + 10 = 24$ centímetro.
\end{solution}

\begin{solution}{exo:g4:measure:4}
Longitud $+$ ancho $= 26 \div 2 = 13$ cm, entonces el ancho es
$13 - 8 = 5$ centímetro.
\end{solution}

\begin{solution}{exo:g4:measure:5}
Muchas formas de L funcionan; por un $10$-cuadrado L hecho de un $4 \times 2$
rectángulo más un $2 \times 1$ pie, caminar la frontera da $16$
unidades. (Cualquier cifra correcta y un conteo cuidadoso es correcto; el
\omterm{def:g4:measure:perimeter}{perímetro} Depende de la L elegida.)
\end{solution}

\begin{solution}{exo:g4:measure:6}
$3 \times 4$: \omterm{def:g4:measure:perimeter}{perímetro} $14$. $2 \times 6$: \omterm{def:g4:measure:perimeter}{perímetro} $16$.
($1 \times 12$: \omterm{def:g4:measure:perimeter}{perímetro} $26$.) Cuanto más cerca de un cuadrado, más
más pequeño el \omterm{def:g4:measure:perimeter}{perímetro} por lo mismo \omterm{def:g4:measure:area}{área}.
\end{solution}

\begin{solution}{exo:g4:measure:7}
\omterm{def:g3:division:half}{Medio} del $6 \times 3$ rectángulo: $18 \div 2 = 9$ cuadrícula.
\end{solution}

\begin{solution}{exo:g4:measure:8}
$10{:}20 \to 12{:}45$: $2$ h $25$ mín. $7{:}35 \to 9{:}10$:
$1$ h $35$ mín. Película: $20{:}30 + 1$ h $50 = 22{:}20$.
\end{solution}

\begin{solution}{exo:g4:measure:9}
Una vuelta: $2 \times (60 + 45) = 210$ metro. Tres vueltas:
$3 \times 210 = 630$ metro.
\end{solution}

\begin{solution}{exo:g4:measure:10}
\omterm{def:g4:measure:perimeter}{perímetro}: $4 \times 85 = 340$ metro; menos la puerta:
$340 - 4 = 336$ m de valla.
\end{solution}

\begin{solution}{exo:g4:measure:11}
Ambas afirmaciones son \emph{FALSO}. Mismo \omterm{def:g4:measure:perimeter}{perímetro},
diferente \omterm{def:g4:measure:area}{áreas}: el
$4 \times 2$ rectángulo y la forma de L de la figura del capítulo (ambos
\omterm{def:g4:measure:perimeter}{perímetro} $12$; \omterm{def:g4:measure:area}{áreas} $8$ y $6$). Mismo \omterm{def:g4:measure:area}{área},
diferente
\omterm{def:g4:measure:perimeter}{perímetros}: el $4 \times 4$ cuadrado y el $8 \times 2$ rectángulo
(ambos \omterm{def:g4:measure:area}{área} $16$; \omterm{def:g4:measure:perimeter}{perímetros} $16$ y $20$).
\end{solution}
